\chapter{Kết luận và Hướng Phát triển}

\section{Tổng kết các Kết quả Đạt được}

Đồ án đã thiết kế, hiện thực hóa và đánh giá có kiểm soát một hệ thống Tác nhân AI Kép cho chuỗi bán lẻ tại Việt Nam gồm hai tác nhân chuyên biệt (\texttt{Retail-AI-} và \texttt{ops\_chest}) kết nối qua cầu nối dữ liệu--API \texttt{ops-retail-bridge}. Đối chiếu tương ứng 1--1 với bốn mục tiêu cụ thể đặt ra ở Chương~1, kết quả nghiên cứu đạt được gồm:

\begin{enumerate}
    \item \textbf{Đạt Mục tiêu 1 -- AI Agent Tư vấn Khách hàng (\emph{Zero-Planner Hybrid Routing}):} Thiết kế cơ chế định tuyến xác định trước -- LLM có giới hạn sau giúp giải quyết 47/50 (94,0\%) câu hỏi trên tập kiểm thử độc lập 50 câu (100\% căn cứ SKU danh mục, 10/10 cổng an toàn/OWASP) với phân vị độ trễ \textbf{P50 = 74,6\,ms, P95 = 139,0\,ms, P99 = 154,5\,ms} (cổng tiền-LLM hạ P99 đi 2,8$\times$), đạt 100/100 trên bộ Bán chạy \& Tồn kho với 100\% minh bạch nguồn \texttt{SIMULATED}/snapshot, đạt 94,0\% trung thực mệnh đề trên tập kiểm toán chuyên gia 100 câu độc lập, đồng thời làm giàu sạch 0 lỗi cho 2.205/3.099 SKU (71,15\%) và huấn luyện mô hình thị giác \texttt{MobileNetV3Large-224} TFLite đạt 92,67\% Top-1 đa góc và 93,8\% Brand@1 / 100\% chặn OOD ở chế độ Hybrid.
    \item \textbf{Đạt Mục tiêu 2 -- AI Agent Dự báo và Vận hành Hàng Tươi sống (\emph{Profit-Gated Router \texttt{B-Gate}}):} Trên 50 cửa hàng độc lập của FreshRetailNet-50K, khôi phục nhu cầu theo giờ giảm WAPE từ 32,35\% xuống 20,53\% (độ lệch $-32{,}35\% \to -0{,}15\%$); LightGBM~v3 (29 đặc trưng, huấn luyện trên 14,91\,M mẫu của 848 cửa hàng tách biệt) hạ WAPE dự báo từ 35,06\% (EWMA) xuống 30,17\% (độ phủ khoảng 80\% đạt 81,2\%). Trong mô phỏng A/B 37 ngày (kiểm định chéo 2 nếp), \texttt{B-Gate} ghép trung vị LightGBM với $\sigma_{\text{EWMA}}$ theo ngưỡng \code{size_cut=30} đạt mức tăng lợi nhuận mô phỏng trung vị \textbf{+4,90\%} so với Baseline~A và \textbf{+0,78\%} so với EWMA (95\% bootstrap CI $[+0{,}48\%;\,+1{,}21\%]$, Wilcoxon $p = 6{,}39\times 10^{-6}$, thắng 76\% cửa hàng; 86,5\% trên nhóm cửa hàng $>30$ SKU), hạ tỷ lệ hủy còn 4,82\% và giảm đứt hàng nhóm hút khách từ 8,69\% xuống 3,37\%, đồng thời loại bỏ hoàn toàn công cụ tự thực thi để bảo đảm phê duyệt HITL.
    \item \textbf{Đạt Mục tiêu 3 -- Phối hợp và Trao đổi Dữ liệu Liên Agent:} Xây dựng bảng ánh xạ dùng chung 164/164 SKU tươi sống (\code{data/bridge/crosswalk.csv}) và lớp cầu nối chỉ đọc fail-closed (\code{backend/services/ops_bridge.py}), chứng minh qua bộ kiểm thử đầu-cuối trên môi trường mô phỏng (\code{tests/test_ops_bridge_e2e.py}) rằng giá giảm FEFO đã duyệt lan truyền tức thì sang Agent khách hàng, đề xuất \texttt{PENDING} bị chặn hoàn toàn, sản phẩm hết hàng bị ẩn khỏi gợi ý còn hàng và hệ thống tự thoái lui an toàn khi cầu nối quá hạn TTL hoặc ngắt kết nối.
    \item \textbf{Đạt Mục tiêu 4 -- Phương pháp Đánh giá Khoa học và Trung thực:} Thiết lập quy trình chống rò rỉ dữ liệu rõ ràng giữa tập phát triển và các tập kiểm thử độc lập, thực hiện ablation study có kiểm soát (A1--A4 và B0--B5), báo cáo khoảng tin cậy, phân tách cỡ cửa hàng, phân tích độ nhạy 8 kịch bản, và định lượng rõ các điểm mù của hậu kiểm số liệu \code{guard.py} cũng như kết quả đảo chiều khử nhiễu ở 5\,dB.
\end{enumerate}

\section{Hạn chế và Hướng Phát triển Tương lai}

Nhằm giữ tính trung thực khoa học trong phạm vi đã kiểm chứng, năm hạn chế chính và hướng nghiên cứu kế tiếp được xác định như sau:
\begin{enumerate}
    \item \textbf{Triển khai Thực địa và Kết nối POS/ERP Thực tế:} Kết nối cầu nối \code{ops_bridge.py} với hệ thống POS/ERP thật của một chuỗi cửa hàng tiện lợi và tiến hành thí điểm A/B 4 tuần trên 3--5 cửa hàng vật lý để đo lường mức tăng lợi nhuận thực, độ trễ giao hàng nhà cung cấp và tỷ lệ chấp nhận đề xuất HITL của quản lý ngoài mô phỏng.
    \item \textbf{Nâng cấp Kiểm soát Số liệu theo Ràng buộc Thực thể (Entity-Bound Grounding):} Mở rộng \code{guard.py} kiểm chứng toàn bộ số lượng, phần trăm giảm giá và số ngày hạn sử dụng dưới ngưỡng 1.000 kèm ánh xạ đúng cặp \emph{(SKU, thuộc tính, giá trị)} nhằm khắc phục điểm mù phát hiện ở Mục~5.5.
    \item \textbf{Mở rộng Tập Kiểm thử Giọng nói Đa Vùng miền:} Thu thập bộ kiểm thử tối thiểu 100--200 câu nói tiếng Việt đa người nói (giọng Bắc--Trung--Nam, tốc độ và khoảng cách micro khác nhau) trong âm vang siêu thị thật để đánh giá thống kê độ tin cậy của ASR và thiết kế ngưỡng bật/tắt khử nhiễu động tránh suy giảm thanh điệu ở SNR thấp ($\le 5$\,dB).
    \item \textbf{Mở rộng Độ phủ Nhãn Dinh dưỡng và Từ điển Đa ngôn ngữ:} Tiếp tục bổ sung nhãn thành phần/calo cho 55,7\% SKU chưa có dữ liệu công khai và chuẩn hóa bộ ánh xạ song ngữ Anh--Việt (khắc phục lỗi truy xuất như \emph{``soy sauce''} trên tập kiểm thử độc lập).
    \item \textbf{Học Chính sách Đặt hàng Đầu-Cuối (End-to-End Decision-Focused Learning):} Nghiên cứu thay thế bước cổng phân khúc rời rạc \texttt{B-Gate} bằng hàm mất mát định hướng quyết định SPO+~\cite{elmachtoub2022spo} hoặc tối ưu tồn kho liên chi nhánh, đồng thời đánh giá thực nghiệm các mô hình ngôn ngữ lượng tử hóa tại chỗ để đo lường chính xác độ trễ và độ chính xác thay vì suy đoán lý thuyết.
\end{enumerate}
